\documentclass[10pt,a4paper]{article}
\usepackage[slovene]{babel}
\usepackage{fontspec}
\usepackage{geometry}
\geometry{left=21mm,right=21mm,top=20mm,bottom=19mm}
\usepackage{microtype}
\usepackage{xcolor}
\usepackage{enumitem}
\usepackage{hyperref}
\usepackage{titlesec}
\usepackage{amsmath}
\definecolor{navy}{HTML}{1D3348}
\hypersetup{colorlinks=true,urlcolor=navy}
\titleformat{\section}{\bfseries\color{navy}\normalsize}{}{0pt}{}
\titlespacing*{\section}{0pt}{8pt}{2pt}
\setlength{\parindent}{0pt}
\setlength{\parskip}{3pt}
\setlist[itemize]{leftmargin=1.35em,itemsep=1pt,topsep=2pt,parsep=0pt}
\pagestyle{empty}

\begin{document}

\begin{center}
    {\fontsize{11}{13}\selectfont Prijava teme za seminar QCA}\par
    \vspace{5pt}
    {\bfseries\fontsize{16}{19}\selectfont Analiza odpornosti QCA demultiplekserja\par na napake zaradi premikov celic}\par
    \vspace{6pt}
    {\small Skupina: 27\qquad Člani: Vojko Hysz, Gašper Suhadolnik, Matic Poženel}
\end{center}

\vspace{-5pt}
\section*{Izhodišče}
Izhajali bomo iz članka Paramjeet in sod. (2026), ki predstavi izdelavo modularnih demultiplekserjev $1:2^n$.Demultiplekser $1:2$ je realiziran s \textbf{27 celicami QCA} in simuliran v QCADesignerju.

Ta članek že navaja, da je za nadaljnje delo med drugim ustrezna tudi analiza odpornosti na različne napake.

\section*{Raziskovalno vprašanje}
Kako premiki posameznih celic vplivajo na pravilnost delovanja, polarizacijo izhodov in energijsko disipacijo demultiplekserja 1:2, zasnovanega s kvantnimi celičnimi avtomati (QCA)? Katere celice so najbolj občutljive na tovrstne napake?

\section*{Hipoteza}
Predpostavljamo, da občutljivost vezja na premike celic ni enakomerna. Premiki celic v bližini logičnih vrat (invertorja in majoritetnih vrat) bodo povzročili večje spremembe izhodne polarizacije in hitreje privedli do nepravilnega delovanja kot primerljivi premiki celic v ravnih povezovalnih vodilih.

Pričakujemo tudi, da bodo večji premiki povzročili izrazitejše odstopanje od rezultatov izvirne implementacije.

\section*{Metrike}
\begin{description}
    \item[Pravilnost delovanja] Stanja vhodnih in izhodnih QCA celic se morajo ujemati z resničnostno tabelo demultiplekserja pri vseh štirih kombinacijah vhodov.
    \item [Izhodna polarizacija] Kolikšni sta velikost in stabilnost polarizacije izhodnih celic v ustreznih fazah urnega cikla?
    \item[Toleranca na premik] Kakšen je največji preizkušeni premik posamezne celice, pri katerem vezje še pravilno deluje?
    \item[Energijska disipacija] Kako se ob premikih spremeni poraba energije glede na originalno vezje?
    \item[Zakasnitev] Kakšna je sprememba števila urinih ciklov, potrebnih za pridobitev pravilnih izhodov?
\end{description}

\section*{Načrt poskusov}

Uporabili bomo 27-celični demultiplekser $1:2$ iz izhodiščnega članka. V simulatorju bomo rekonstruirali izvirno vezje ter preverili njegovo pravilnost in rezultate primerjali z objavljenimi meritvami.

Nato bomo izbrali štiri reprezentativne celice na različnih mestih vezja: v NOT vratih, v majoritetnih vratih, na vhodnem vodilu in na izhodnem vodilu. Vsako bomo ločeno premaknili za $1$, $3$ in $5$ nm v določeni smeri, pri čemer bodo preostale celice ostale nespremenjene.

Tako bomo pripravili 12 spremenjenih različic in izvirno vezje. Za vsako različico bomo izvedli simulacijo, preverili pravilnost izhodov ter izmerili izbrane metrike.

Rezultate bomo ustrezno prikazali in ugotovili, katere vrste celic so najbolj občutljive na premike. Če se pri izbranih premikih napake ne pojavijo, bomo območje premikov ustrezno razširili.

\section*{Orodja in razdelitev dela}

\textbf{Orodja:} QCADesigner-E (zasnova, funkcijska in energijska simulacija), Python (obdelava meritev in grafi), \LaTeX{} (poročilo).

\begin{itemize}
    \item \textbf{1. član:} Rekonstrukcija izvirnega demultiplekserja, preverjanje logične pravilnosti in primerjava z referenčnim člankom.
    \item \textbf{2. član:} Priprava različic vezja s premaknjenimi celicami in izvedba simulacij.
    \item \textbf{3. član:} Analiza polarizacije, zakasnitev in energijske disipacije, statistična obdelava rezultatov ter priprava grafov.
\end{itemize}

\vspace{3pt}
{\footnotesize\textbf{Izhodiščni članek:} Paramjeet in sod., \emph{Low Complexity QCA-Based 1:2\textsuperscript{n} Demux Architecture With Effective Modular Design}, Journal of Engineering, 2026. DOI: \href{https://doi.org/10.1155/je/9541510}{10.1155/je/9541510}.}
\end{document}
